\begin{lemma}[Transfer of neighborhood bounds to a codegree window]
Let $\varepsilon,p,q\in[0,1]$, and let $k,a,c$ be nonnegative integers with $a\leq k$.
Suppose
\[
 (p-\varepsilon)_+k\leq a\leq(p+\varepsilon)k,
 \qquad(q-\varepsilon)_+a\leq c\leq(q+\varepsilon)a.
\]
Then $(pq-2\varepsilon)_+k\leq c\leq(pq+2\varepsilon)k$.
The subscript $+$ denotes truncated subtraction, $(x-y)_+=\max\{x-y,0\}$.
\end{lemma}
\begin{proof}
Multiplication of a truncated difference by a finite nonnegative integer distributes over truncated subtraction.
Thus the two lower bounds give $pk\leq a+\varepsilon k$ and $qa\leq c+\varepsilon a$.
Since $q\leq1$ and $a\leq k$, monotonicity and distributivity give
\[
 pqk=q(pk)\leq q(a+\varepsilon k)
 =qa+q\varepsilon k\leq c+\varepsilon a+\varepsilon k
 \leq c+2\varepsilon k.
\]
The defining order property of truncated subtraction yields $(pqk-2\varepsilon k)_+\leq c$; distributing the finite factor $k$ gives the required lower bound.
For the upper bound,
\[
 c\leq(q+\varepsilon)a=qa+\varepsilon a
 \leq q(p+\varepsilon)k+\varepsilon k
 =pqk+q\varepsilon k+\varepsilon k
 \leq(pq+2\varepsilon)k.
\]
No cancellation or division by a parameter occurs, so the argument includes zero values of every parameter and of $k$.
\end{proof}
